\chapter{Непрерывные функции}

% -----------------------------------------------------------------------------
\lecture{Основные определения и примеры}

\subsection{Непрерывность функции в точке}

\begin{definition}{Непрерывность функции в точке}{c4-continuity-point-initial}
  Функция $f$ называется непрерывной в точке $a$, если для любой
  окрестности $V(f(a))$ значения $f(a)$ функции в точке $a$ найдется такая
  окрестность $U(a)$ точки $a$, образ которой при отображении $f$
  содержится в $V(f(a))$.

  Приведем формально-логическую запись этого определения вместе с двумя
  его вариациями, часто используемыми в анализе:
  \[
    (f\text{ непрерывна в точке }a)
    \;:=\;
    \forall V(f(a))\ \exists U(a)\
    \bigl(f(U(a))\subset V(f(a))\bigr),
  \]
  \[
    \forall\varepsilon>0\ \exists U(a)\ \forall x\in U(a)\
    (|f(x)-f(a)|<\varepsilon),
  \]
  \[
    \forall\varepsilon>0\ \exists\delta>0\ \forall x\in\R\
    \bigl(|x-a|<\delta\Rightarrow|f(x)-f(a)|<\varepsilon\bigr).
  \]
\end{definition}

\begin{definition}{Непрерывность функции на произвольном множестве в точке}{c4-continuity-point}
  Функция $f\colon E\to\R$ называется непрерывной в точке $a\in E$, если
  для любой окрестности $V(f(a))$ значения $f(a)$ функции, принимаемого ею
  в точке $a$, найдется такая окрестность $U_E(a)$ точки $a$ в множестве
  $E$, образ которой $f(U_E(a))$ содержится в $V(f(a))$.

  Итак,
  \[
    (f\colon E\to\R\text{ непрерывна в }a\in E)
    \;:=\;
    \forall V(f(a))\ \exists U_E(a)\
    \bigl(f(U_E(a))\subset V(f(a))\bigr).
  \]
  Равносильные $\varepsilon$-формы определения:
  \[
    \forall\varepsilon>0\ \exists U_E(a)\ \forall x\in U_E(a)\
    (|f(x)-f(a)|<\varepsilon),
  \]
  \[
    \forall\varepsilon>0\ \exists\delta>0\ \forall x\in E\
    \bigl(|x-a|<\delta\Rightarrow|f(x)-f(a)|<\varepsilon\bigr).
  \]
\end{definition}

\begin{properties}{Непрерывность функции в точке}{c4-continuity-point-properties}
  \begin{enumerate}[label=\arabic*${}^{\circ}$.]
    \item В любой изолированной точке области определения функция,
      очевидно, непрерывна.
    \item Если $a\in E$ и $a$ \textemdash{} предельная точка множества $E$,
      то
      \[
        (f\colon E\to\R\text{ непрерывна в }a)
        \quad\Longleftrightarrow\quad
        \lim_{E\ni x\to a}f(x)=f(a).
      \]
    \item Непрерывные в точке функции (операции) и только они
      перестановочны с операцией предельного перехода:
      \[
        \lim_{E\ni x\to a}f(x)
        =f\!\left(\lim_{E\ni x\to a}x\right).
      \]
    \item Если при $a\in E$ окрестности $U_E(a)$ точки $a$ образуют базу
      $\mathcal B_a$, то
      \[
        (f\colon E\to\R\text{ непрерывна в }a\in E)
        \quad\Longleftrightarrow\quad
        \lim_{\mathcal B_a}f(x)=f(a).
      \]
    \item Непрерывность функции $f\colon E\to\R$ в точке $a\in E$
      равносильна существованию предела этой функции по базе
      $\mathcal B_a$ окрестностей $U_E(a)$ точки $a$ в $E$:
      \[
        (f\colon E\to\R\text{ непрерывна в }a\in E)
        \quad\Longleftrightarrow\quad
        \exists\lim_{\mathcal B_a}f(x).
      \]
    \item Функция непрерывна в точке $a\in E$ тогда и только тогда, когда
      для любого $\varepsilon>0$ найдется окрестность $U_E(a)$ точки $a$ в
      $E$ такая, на которой колебание $\omega(f;U_E(a))$ функции меньше
      $\varepsilon$.
  \end{enumerate}
\end{properties}

\begin{definition}{Колебание функции в точке}{c4-oscillation-at-point}
  Величина
  \[
    \omega(f;a)=\lim_{\delta\to+0}\omega(f;U_E^\delta(a)),
  \]
  где $U_E^\delta(a)$ есть $\delta$-окрестность точки $a$ в множестве $E$,
  называется колебанием функции $f\colon E\to\R$ в точке $a$.
\end{definition}

\begin{properties}{Критерий непрерывности через колебание}{c4-oscillation-continuity}
  Функция непрерывна в точке тогда и только тогда, когда ее колебание в
  этой точке равно нулю:
  \[
    (f\colon E\to\R\text{ непрерывна в }a\in E)
    \quad\Longleftrightarrow\quad
    \omega(f;a)=0.
  \]
\end{properties}

\begin{definition}{Непрерывность на множестве}{c4-continuity-on-set}
  Функция $f\colon E\to\R$ называется непрерывной на множестве $E$, если
  она непрерывна в каждой точке множества $E$.

  Совокупность всех вещественнозначных функций, непрерывных на множестве
  $E$, обозначается символом $C(E;\R)$ или, короче, $C(E)$.
\end{definition}

\subsection{Точки разрыва}

\begin{definition}{Точка разрыва}{c4-discontinuity-point}
  Если функция $f\colon E\to\R$ не является непрерывной в некоторой точке
  множества $E$, то эта точка называется точкой разрыва функции $f$.

  Формально:
  \[
    (a\in E\text{ \textemdash{} точка разрыва функции }f)
    \;:=\;
    \exists V(f(a))\ \forall U_E(a)\ \exists x\in U_E(a)\
    \bigl(f(x)\notin V(f(a))\bigr).
  \]
  В $\varepsilon$-$\delta$-форме:
  \[
    \exists\varepsilon>0\ \forall\delta>0\ \exists x\in E\
    \bigl(|x-a|<\delta\ \land\ |f(x)-f(a)|>\varepsilon\bigr).
  \]
\end{definition}

\begin{definition}{Точка устранимого разрыва}{c4-removable-discontinuity}
  Если точка разрыва $a\in E$ функции $f\colon E\to\R$ такова, что
  существует непрерывная функция $\widetilde f\colon E\to\R$ такая, что
  \[
    f|_{E\setminus\{a\}}=\widetilde f|_{E\setminus\{a\}},
  \]
  то $a$ называется точкой устранимого разрыва функции $f\colon E\to\R$.

  Точка устранимого разрыва характеризуется тем, что существует предел
  \[
    \lim_{E\ni x\to a}f(x)=A,
  \]
  но $A\ne f(a)$, и достаточно положить
  \[
    \widetilde f(x)=
    \begin{cases}
      f(x), & x\in E,\ x\ne a,\\
      A, & x=a,
    \end{cases}
  \]
  как мы уже получим непрерывную в точке $a$ функцию
  $\widetilde f\colon E\to\R$.
\end{definition}

\begin{definition}{Точка разрыва первого рода}{c4-first-kind-discontinuity}
  Точка $a\in E$ называется точкой разрыва первого рода для функции
  $f\colon E\to\R$, если существуют пределы
  \[
    \lim_{E\ni x\to a-0}f(x)=:f(a-0),
    \qquad
    \lim_{E\ni x\to a+0}f(x)=:f(a+0),
  \]
  но по крайней мере один из этих пределов не совпадает со значением $f(a)$
  функции в точке $a$.

  Если все точки множества $E$ в некоторой окрестности точки $a$ лежат по
  одну сторону от точки $a$, рассматривается только один из указанных
  пределов.
\end{definition}

\begin{definition}{Точка разрыва второго рода}{c4-second-kind-discontinuity}
  Если $a\in E$ \textemdash{} точка разрыва функции $f\colon E\to\R$ и в
  этой точке не существует по меньшей мере один из пределов, указанных в
  определении точки разрыва первого рода, то $a$ называется точкой разрыва
  второго рода.

  Таким образом, имеется в виду, что всякая точка разрыва, не являющаяся
  точкой разрыва первого рода, является точкой разрыва второго рода.
\end{definition}

% -----------------------------------------------------------------------------
\lecture{Свойства непрерывных функций}

\subsection{Локальные свойства}

\begin{theorem}{Локальные свойства непрерывных функций}{c4-local-properties}
  Пусть $f\colon E\to\R$ \textemdash{} функция, непрерывная в точке
  $a\in E$. Тогда справедливы следующие утверждения:
  \begin{enumerate}[label=\arabic*${}^{\circ}$.]
    \item Функция $f$ ограничена в некоторой окрестности $U_E(a)$ точки
      $a$.
    \item Если $f(a)\ne0$, то в некоторой окрестности $U_E(a)$ точки $a$
      все значения функции положительны или отрицательны вместе с $f(a)$.
    \item Если функция $g\colon U_E(a)\to\R$ определена в некоторой
      окрестности точки $a$ и, как и $f\colon E\to\R$, непрерывна в самой
      точке $a$, то функции:
      \begin{enumerate}[label=\alph*)]
        \item $(f+g)(x):=f(x)+g(x)$,
        \item $(f\cdot g)(x):=f(x)\cdot g(x)$,
        \item $\left(\dfrac fg\right)(x):=\dfrac{f(x)}{g(x)}$
          (при условии, что $g(x)\ne0$)
      \end{enumerate}
      определены в некоторой окрестности точки $a$ и непрерывны в точке
      $a$.
    \item Если функция $g\colon Y\to\R$ непрерывна в точке $b\in Y$, а
      функция $f$ такова, что $f\colon E\to Y$, $f(a)=b$ и $f$ непрерывна в
      точке $a$, то композиция $g\circ f$ определена на $E$ и также
      непрерывна в точке $a$.
  \end{enumerate}
\end{theorem}

\subsection{Глобальные свойства непрерывных функций}

\begin{theorem}{Больцано\textendash{}Коши о промежуточном значении}{c4-bolzano-cauchy}
  Если функция, непрерывная на отрезке, принимает на его концах значения
  разных знаков, то на отрезке есть точка, в которой функция обращается в
  нуль.

  В логической символике:
  \[
    (f\in C[a,b])\ \land\ (f(a)f(b)<0)
    \quad\Longrightarrow\quad
    \exists c\in[a,b]\;(f(c)=0).
  \]
\end{theorem}

\begin{corollary}{Теоремы Больцано\textendash{}Коши}{c4-intermediate-value-corollary}
  Если функция $\varphi$ непрерывна на интервале и в каких-то точках $a$ и
  $b$ интервала принимает значения $\varphi(a)=A$ и $\varphi(b)=B$, то для
  любого числа $C$, лежащего между $A$ и $B$, найдется точка $c$, лежащая
  между точками $a$ и $b$, в которой $\varphi(c)=C$.
\end{corollary}

\begin{theorem}{Вейерштрасса о максимальном значении}{c4-weierstrass-extreme-value}
  Функция, непрерывная на отрезке, ограничена на нем. При этом на отрезке
  есть точка, где функция принимает максимальное значение, и есть точка,
  где она принимает минимальное значение.
\end{theorem}

\begin{definition}{Равномерная непрерывность}{c4-uniform-continuity}
  Функция $f\colon E\to\R$ называется равномерно непрерывной на множестве
  $E\subset\R$, если для любого числа $\varepsilon>0$ найдется число
  $\delta>0$ такое, что для любых точек $x_1,x_2\in E$ таких, что
  $|x_1-x_2|<\delta$, выполнено $|f(x_1)-f(x_2)|<\varepsilon$.

  Короче,
  \[
    (f\colon E\to\R\text{ равномерно непрерывна})
    \;:=\;
    \forall\varepsilon>0\ \exists\delta>0\ \forall x_1\in E\ \forall x_2\in E
  \]
  \[
    \bigl(|x_1-x_2|<\delta
    \Rightarrow |f(x_1)-f(x_2)|<\varepsilon\bigr).
  \]
  \par
\end{definition}

\begin{properties}{Равномерная непрерывность}{c4-uniform-continuity-properties}
  \begin{enumerate}[label=\arabic*${}^{\circ}$.]
    \item Если функция равномерно непрерывна на множестве, то она
      непрерывна в любой его точке.
    \item Непрерывность функции, вообще говоря, не влечет ее равномерную
      непрерывность.
  \end{enumerate}

  Отрицание свойства функции быть равномерно непрерывной имеет вид
  \[
    (f\colon E\to\R\text{ не является равномерно непрерывной})
    \;:=\;
    \exists\varepsilon>0\ \forall\delta>0\ \exists x_1\in E\ \exists x_2\in E
  \]
  \[
    \bigl(|x_1-x_2|<\delta
    \ \land\ |f(x_1)-f(x_2)|\ge\varepsilon\bigr).
  \]
  \par
\end{properties}

\begin{theorem}{Кантора\textendash{}Гейне о равномерной непрерывности}{c4-cantor-heine}
  Функция, непрерывная на отрезке, равномерно непрерывна на этом отрезке.
\end{theorem}

\begin{proposition}{Инъективность непрерывной функции на отрезке}{c4-injective-monotone}
  Непрерывное отображение $f\colon E\to\R$ отрезка $E=[a,b]$ в $\R$
  инъективно в том и только в том случае, когда функция $f$ строго
  монотонна на отрезке $[a,b]$.
\end{proposition}

\begin{proposition}{Обратная к строго монотонной функции}{c4-inverse-monotone}
  Каждая строго монотонная функция $f\colon X\to\R$, определенная на
  числовом множестве $X\subset\R$, обладает обратной функцией
  $f^{-1}\colon Y\to\R$, которая определена на множестве $Y=f(X)$ значений
  функции $f$ и имеет на $Y$ тот же характер монотонности, какой имеет
  функция $f$ на множестве $X$.
\end{proposition}

\begin{proposition}{Разрывы монотонной функции}{c4-monotone-discontinuities}
  Функция $f\colon E\to\R$, монотонная на множестве $E\subset\R$, может
  иметь на $E$ разрывы только первого рода.
\end{proposition}

\begin{corollary}{Структура разрыва монотонной функции}{c4-monotone-discontinuity-structure}
  Если $a$ \textemdash{} точка разрыва монотонной функции
  $f\colon E\to\R$, то по крайней мере один из пределов
  \[
    \lim_{E\ni x\to a-0}f(x)=f(a-0),
    \qquad
    \lim_{E\ni x\to a+0}f(x)=f(a+0)
  \]
  определен; по крайней мере в одном из неравенств
  $f(a-0)\le f(a)\le f(a+0)$, если $f$ \textemdash{} неубывающая (или
  $f(a-0)\ge f(a)\ge f(a+0)$, если $f$ \textemdash{} невозрастающая)
  функция, имеет место знак строгого неравенства; в интервале, определяемом
  этим строгим неравенством, нет ни одного значения функции; указанные
  интервалы, отвечающие различным точкам разрыва монотонной функции, не
  пересекаются.
\end{corollary}

\begin{corollary}{Множество точек разрыва монотонной функции}{c4-countable-discontinuities}
  Множество точек разрыва монотонной функции не более чем счетно.
\end{corollary}

\begin{proposition}{Критерий непрерывности монотонной функции}{c4-monotone-continuity-criterion}
  Монотонная функция $f\colon E\to\R$, заданная на отрезке $E=[a,b]$,
  непрерывна на нем тогда и только тогда, когда множество $f(E)$ ее значений
  само является отрезком с концами $f(a)$ и $f(b)$.
\end{proposition}

\begin{theorem}{Об обратной функции}{c4-inverse-function}
  Функция $f\colon X\to\R$, строго монотонная на множестве $X\subset\R$,
  имеет обратную функцию $f^{-1}\colon Y\to\R$, определенную на множестве
  $Y=f(X)$ значений функции $f$. Функция $f^{-1}\colon Y\to\R$ монотонна и
  имеет на $Y$ тот же вид монотонности, какой имеет функция
  $f\colon X\to\R$ на множестве $X$.

  Если, кроме того, $X$ есть отрезок $[a,b]$ и функция $f$ непрерывна на
  нем, то множество $Y=f(X)$ есть отрезок с концами $f(a)$, $f(b)$ и функция
  $f^{-1}\colon Y\to\R$ непрерывна на нем.
\end{theorem}
